% Section commune : analyse critique transversale
\section{Analyse critique : ce que les chiffres publiés permettent de conclure}
\label{sec:critique}

\subsection{La puissance : rarement mesurée sur la carte, presque jamais sur la puce FPGA}

Recompte de \texttt{benchmarks.csv} sur les 216 lignes mesurées sur FPGA (44 articles) :

\begin{table}[!htbp]
\centering\small
\caption{Portée de la puissance dans les lignes FPGA de \texttt{benchmarks.csv}
(01/10/2026).}
\label{tab:portee}
\begin{tabularx}{\textwidth}{@{}lrL@{}}
\toprule
Portée & Lignes & Origine \\
\midrule
\emph{board} (carte au wattmètre) & 46 & 9 articles : Zhang 2024 10, FINN 8, FINN-R 8,
  Elastic-DF 8, \glslink{lp-yolo}{LP-YOLO} 6, Schlosser 3, Zhang 2015, \glslink{onlad}{ONLAD} et \glslink{codenet}{CoDeNet} 1 chacun \\
\emph{chip} & 1 & Castelino (rails PMBus de la \glslink{zcu104}{ZCU104}) \\
\emph{board-dynamic} & 0 & --- \\
\emph{estimated} (outil) & 16 & \glslink{ef-train}{EF-Train} 6, Manjunath 4, \glslink{imagehd}{ImageHD} 3, Karki 2, Que 2024 1 \\
\emph{unspecified} & 52 & dont SATAY 14, Piyasena 16 (deux articles), Zhai 8 \\
puissance non rapportée (\nr) & 101 & --- \\
\bottomrule
\end{tabularx}
\end{table}

Un peu plus de la moitié des lignes FPGA (115 sur 216) ont une puissance, et 47
seulement une puissance mesurée, issues de 10 articles. Les 18 lignes \emph{chip} de la base viennent,
sauf une, de \glslink{orin}{Jetson Orin} (12, \citealp{arxiv-2502.15737-performance-analysis-you-only-look}),
du Raspberry~Pi~5 (2) et de l'ASIC \glslink{vega}{VEGA} (3). La seule ligne \emph{chip} FPGA est celle de
\citet[§IV-B]{arxiv-2407.17647-energy-efficient-artefact-detection-accelerator} : 5,38\,\glslink{u-w}{W}
actifs pour 4,6\,\glslink{u-w}{W} au repos. Les lignes \emph{board} ne sont plus surtout anciennes : à
côté de FINN, FINN-R, Zhang 2015 et \glslink{onlad}{ONLAD}, elles viennent de
\citet{2024-zhang-end-end-acceleration-yolo-object},
\citet{2021-alonso-elastic-df-scaling-performance-dnn} (deux Alveo, hors embarqué),
\citet{arxiv-2304.11580-framework-benchmarking-real-time-embedded} et
\citet{arxiv-2609.24757-npu-accelerator-quantized-real-time}\pp. Une seule paire puce /
carte existe sur une même FPGA, et la puce y est estimée : 7,8\,\glslink{u-w}{W} par Vivado contre
12,3\,\glslink{u-w}{W} au wattmètre sur KC705 \citep[Table~3]{2024-zhang-end-end-acceleration-yolo-object}.
Seuls Schlosser et al. donnent aussi la puissance relative, repos soustrait, que le CSV
ne reprend pas faute de ligne \emph{board-dynamic} : 7,0\,\glslink{u-w}{W} sur \glslink{zcu104}{ZCU104} pour 23\,\glslink{u-w}{W} absolus
\src{Table~1}.

Plusieurs comparaisons publiées mettent face à face des portées différentes :
\begin{itemize}
  \item une puissance dynamique de quelques mW contre des puissances totales
    \citep[Table~16]{2019-coutinho-deep-neural-network-hardware-implementation} ;
  \item la carte entière contre le CPU seul et le GPU seul
    \citep[§6.4]{2019-tsukada-neural-network-device-learning-anomaly} ;
  \item le TDP d'un Xeon pris comme puissance mesurée
    \citep[§5.2]{2015-zhang-optimizing-fpga-accelerator-design-deep} ;
  \item une estimation de l'outil contre des mesures par nvidia-smi
    \citep[§VII-D]{arxiv-2604.21280-imagehd-energy-efficient-device-continual} ;
  \item les rails PMBus de la puce contre jtop sur une \glslink{xavier-nx}{Xavier NX} bridée à 10\,\glslink{u-w}{W}
    \citep[Table~VI]{arxiv-2407.17647-energy-efficient-artefact-detection-accelerator} ;
  \item une puissance estimée par Vivado présentée comme \enquote{3x moins} que celle de
    SATAY \citep[Résumé, Table~III]{arxiv-2607.13106-no-attention-no-problem-dpu}\pp.
\end{itemize}
Le survey de \citet{2023-shuvo-efficient-acceleration-deep-learning-inference} rappelle,
dans sa partie sur les conditions de test, qu'un processeur peut afficher une très faible
consommation en omettant les accès à la mémoire externe, et qu'il faut séparer résultats
simulés et mesurés \src{§IX-I}.

\subsection{Des métriques qui flattent}

\begin{itemize}
  \item \textbf{Latence amortie} : un débit inverse présenté comme une latence
    \citep[§IV-E]{2024-que-low-latency-variational-autoencoder-fpgas}.
  \item \textbf{GOp/s par \glslink{dsp}{DSP}} : cette métrique efface les écarts de précision (de W1A8 à
    W32A32) et de carte
    \citep[Table~III]{2023-montgomerie-corcoran-satay-streaming-architecture-toolflow-accelerating}.
    Venieris et al. la jugeaient valable seulement dans une même famille de FPGA
    \citep[§3.1]{2018-venieris-toolflows-mapping-convolutional-neural-networks}. Le
    \enquote{+29\,\% face à Vitis AI} de \citet{2024-yan-fpga-yolov5-accelerator-real-time}
    est un \glslink{gopsdsp}{GOPS/DSP} face à un \glslink{yolov3}{YOLOv3} d'un autre travail, pas une mesure de Vitis AI sur la
    même carte \src{Table~2}.
  \item \textbf{Gains de surface face à des designs plus rapides} : \enquote{$-26$\,\% de
    \glslink{ff}{FF}, $-15$\,\% de \glslink{dsp}{DSP}} face à des accélérateurs de 10,5 à 62,9\,\glslink{gops}{GOPS}, pour 2,33\,\glslink{gops}{GOPS} et
    un autre modèle \citep[Table~9]{2025-kim-design-implementation-yolov2-accelerator-zynq}.
  \item \textbf{\glslink{gopsw}{GOPS/W} dynamiques} face à des \glslink{gopsw}{GOPS/W} totaux
    \citep[Table~8]{2023-zhai-fpga-vehicle-detection-tracking-accelerator}.
  \item \textbf{\glslink{gops}{GOPS} projetés} sur un FPGA hypothétique, comparés à des \glslink{gops}{GOPS} mesurés
    \citep[Tables~16--17]{2019-coutinho-deep-neural-network-hardware-implementation}.
  \item \textbf{GOp/s seul}, sans \glslink{fps}{FPS}, pour un détecteur
    \citep[Table~5]{arxiv-1809.04570-finn-r-end-end-deep} ; \enquote{10,65x} en GOp/s
    entre des UNet de tailles différentes, à \glslink{fps}{FPS} voisins
    \citep[§V-D]{arxiv-2403.18921-smof-streaming-modern-cnns-fpgas}\pp.
  \item \textbf{Gains face à un CPU ou un GPU de bureau}, de portée de puissance non dite :
    \enquote{70x / 37x}, alors que le GPU fait 44,64\,\glslink{fps}{FPS} contre 24,3 sur la \glslink{kv260}{KV260}
    \citep[Table~10]{2026-fata-low-cost-fpga-enhanced-cnn} ; gains d'énergie de 258 à
    383 fois qui mêlent une estimation FPGA et des mesures CPU et GPU
    \citep[Table~VII]{arxiv-2604.21280-imagehd-energy-efficient-device-continual}.
  \item \textbf{Conclusion contredite par son propre tableau} : 2,6\,\glslink{fpsw}{FPS/W} présentés comme
    efficaces, moins que quatre des cinq travaux cités
    \citep[§IV-B, Table~I]{arxiv-2507.18174-real-time-object-detection-classification}.
  \item \textbf{Chiffres incohérents entre résumé et texte} : 35x contre 58x
    \citep[§I-B, §IV-C]{arxiv-2402.09780-tinycl-efficient-hardware-architecture-continual}\pp ;
    perte de précision de $\approx$\,5\,\% contre $\approx$\,7\,\%
    \citep[Résumé, §IV-B]{arxiv-2607.13106-no-attention-no-problem-dpu}\pp.
  \item \textbf{Précisions non comparables} : \glslink{iou}{IoU} d'un seul objet (\glslink{dac-sdc}{DAC-SDC}, \glslink{skynet}{SkyNet}\pp), \glslink{ap}{AP}
    à une classe (Zhai), \enquote{accuracy} sur 72 images (Al Amin), \enquote{accuracy}
    non définie d'un débruiteur (arxiv-2301.07050\pp, seule ligne en confiance \emph{low}),
    \glslink{map}{mAP} sur des véhicules filtrés de VOC (\glslink{lp-yolo}{LP-YOLO}\pp), \glslink{map}{mAP} sur un jeu privé à une cible
    \citep[Table~4]{arxiv-2502.15737-performance-analysis-you-only-look}.
\end{itemize}
Le recompte confirme cette dispersion. Les 87 lignes FPGA de détection (17 articles)
emploient 14 libellés de précision différents, et 12 n'ont aucune précision. Seules 22
portent une précision sur COCO, dont 14 pour SATAY, quantification simulée par ONNX
Runtime ; six d'entre elles seulement (Schlosser, Karki\pp) sont mesurées sur la carte de
façon établie. Sur les 216 lignes FPGA, le \glslink{fps}{FPS} manque 132 fois, la latence 88 fois, la
puissance 101 fois et la précision 48 fois. Shuvo et al. le formulent ainsi : \enquote{EI requires a universal and
comprehensive set of quality metrics to execute meaningful comparisons of DL models,
optimization algorithms, and hardware platforms} \src{§XIV-C}. La grille de Venieris et al.
(débit, latence à batch~1, ressources, puissance, précision, mémoire) existe depuis 2018 ;
aucune fiche du corpus ne la reprend en entier.

\subsection{La précision est rarement remesurée sur la cible}

Quantifier et compiler un réseau change sa précision, et la plupart des articles ne la
remesurent pas sur la carte :
\begin{itemize}
  \item \glslink{ap}{AP} identique en FP32 et en INT16 à trois décimales
    \citep[Table~8]{2023-zhai-fpga-vehicle-detection-tracking-accelerator} ;
  \item quantification simulée par ONNX Runtime
    \citep[§IV-A]{2023-montgomerie-corcoran-satay-streaming-architecture-toolflow-accelerating} ;
  \item comparaison au C++ généré, pas au FPGA
    \citep[§8]{arxiv-2101.05108-fast-convolutional-neural-networks-fpgas} ;
  \item \glslink{auroc}{AUC} mesurées en TensorFlow
    \citep[§5]{2019-tsukada-neural-network-device-learning-anomaly} ;
  \item seules des pertes relatives de \glslink{map}{mAP} sont données
    \citep[§6.2]{2019-ding-req-yolo-resource-aware-efficient} ;
  \item $-1,5$\,point de \glslink{map}{mAP} VOC, sans dire si la mesure est faite sur la carte
    \citep[Table~2]{2024-yan-fpga-yolov5-accelerator-real-time} ;
  \item écart entre le modèle QAT et le modèle flottant non rapporté
    \citep{arxiv-2609.24757-npu-accelerator-quantized-real-time}\pp ;
  \item \glslink{map}{mAP} identiques sur \glslink{jetson-nano}{Jetson Nano} et \glslink{orin-nano}{Orin Nano}, sans preuve d'une mesure par carte
    \citep[§IV-C]{arxiv-2409.16808-comprehensive-evaluation-deep-learning-object}\pp.
\end{itemize}
Quand la mesure sur cible est faite, elle révèle une perte, parfois forte et imprévisible
(table~\ref{tab:precision-cible}).

\begin{table}[!htbp]
\centering\footnotesize
\caption{Pertes de précision mesurées après quantification ou compilation. La mesure est
faite sur la cible chez Schlosser, Karki, Rey, Feng, Malviya et Coutinho ; pour Castelino,
SkyNet et SMOF, le texte extrait ne le précise pas.}
\label{tab:precision-cible}
\begin{tabularx}{\textwidth}{@{}P{3.1cm}P{2.6cm}LP{2.2cm}@{}}
\toprule
Travail & Modèle, jeu & Avant $\rightarrow$ après & Source \\
\midrule
Schlosser & \glslink{yolov4}{YOLOv4}, COCO & Vitis AI Quantizer : $-2,2$\,\% de \glslink{map}{mAP}, la compilation \glslink{dpu}{DPU} en
  retire encore ; TensorRT INT8 38,6 contre \glslink{dpu}{DPU} INT8 37,5 (\glslink{map5095}{mAP@0.5:0.95}) & §4.3, Table~1 \\
Karki\pp & \glslink{yolov8n}{YOLOv8n}, KITTI / VOC & 0,64 $\rightarrow$ 0,34 ; 0,60 $\rightarrow$ 0,45
  (CPU contre FPGA INT8) & Table~II \\
Karki\pp & \glslink{yolov11n}{YOLOv11n}, VOC / COCO / personnes & 0,59 $\rightarrow$ 0,57 ; 0,31
  $\rightarrow$ 0,29 ; 0,75 $\rightarrow$ 0,71 & Table~II \\
Rey (\glslink{jetson}{Jetson}, hors FPGA) & \glslink{yolov8n}{YOLOv8n} / v8s, jeu privé & TensorRT INT8 : 84,2 $\rightarrow$
  71,9 ; 86,1 $\rightarrow$ 79,7 (\glslink{map5095}{mAP50-95}) & Table~4 \\
Feng (\glslink{ncs2}{NCS2}, hors FPGA) & \glslink{yolov3-tiny}{YOLOv3-tiny}, vidéo & OpenVINO : confiance moyenne 0\,\% contre
  59,7\,\% sur \glslink{jetson}{Jetson} & Table~3 \\
Castelino & \glslink{cae}{CAE} d'anomalies & INT8 en PTQ seule : \glslink{f1}{F1} 0 ; après fine-tuning 92,77 (FP32
  94,05) & Table~III \\
\glslink{skynet}{SkyNet}\pp & \glslink{skynet}{SkyNet}, \glslink{dac-sdc}{DAC-SDC} & \glslink{iou}{IoU} 0,741 $\rightarrow$ 0,727 (FM9/W11) $\rightarrow$
  0,680 (FM8/W10) & Table~7 \\
SMOF\pp & \glslink{yolov8n}{YOLOv8n}, COCO & 37,10 $\rightarrow$ 35,10 (bfp8) & Table~III \\
Malviya & \glslink{cae}{CAE} de vibrations & 88,8\,\% sur CPU $\rightarrow$ 85,02\,\% sur FPGA & §IV-B \\
Coutinho & \glslink{ssae}{SSAE}, MNIST & 93,4\,\% $\rightarrow$ 93,38\,\% & §V-D1 \\
\bottomrule
\end{tabularx}
\end{table}

Sources : \citet{arxiv-2304.11580-framework-benchmarking-real-time-embedded},
\citet{arxiv-2607.13106-no-attention-no-problem-dpu}\pp,
\citet{arxiv-2502.15737-performance-analysis-you-only-look},
\citet{2022-feng-benchmark-analysis-yolo-performance-edge},
\citet{arxiv-2407.17647-energy-efficient-artefact-detection-accelerator},
\citet{arxiv-1909.09709-skynet-hardware-efficient-method-object}\pp,
\citet{arxiv-2403.18921-smof-streaming-modern-cnns-fpgas}\pp,
\citet{2022-malviya-edge-compatible-convolutional-autoencoder-implemented} et
\citet{2019-coutinho-deep-neural-network-hardware-implementation}. Trois constats :
\begin{itemize}
  \item les deux seuls détecteurs mesurés sur une carte FPGA (Schlosser, Karki\pp) perdent
    jusqu'à 30 points de \glslink{map5095}{mAP@0.5:0.95} selon le modèle et le jeu. Chez Karki et al., les
    écarts ne s'expliquent pas : \glslink{yolov11n}{YOLOv11n} perd 2 points sur VOC, \glslink{yolov8n}{YOLOv8n} en perd 15 sur VOC
    et 30 sur KITTI. Les modifications pour le \glslink{dpu}{DPU} (\glslink{relu}{ReLU}, attention approchée) coûtent
    en plus $\approx$\,3 points en moyenne \src{§IV-B, Table~I} ;
  \item les auteurs minimisent parfois la perte qu'ils mesurent : Rey et al. concluent
    \enquote{without a significant loss of accuracy} \src{§6} malgré 12,3 points perdus
    par \glslink{yolov8n}{YOLOv8n} en INT8 ;
  \item la perte dépend de la chaîne autant que du réseau : la PTQ seule effondre
    l'\glslink{ae}{autoencodeur} de Castelino et al., et TensorRT garde une meilleure \glslink{map}{mAP} que Vitis AI
    sur le même \glslink{yolov4}{YOLOv4}.
\end{itemize}
Pour un détecteur de personnes certifiable, cet écart entre le modèle évalué et le modèle
déployé est aussi un problème de preuve \deduit\ifcsname r@sec:surete\endcsname{} (voir la
section~\ref{sec:surete})\fi.

\subsection{Baselines GPU embarqué faibles ou absentes}

Quand un article FPGA se compare à un GPU embarqué, la baseline est souvent ancienne,
bridée ou sous-optimisée :
\begin{itemize}
  \item \glslink{tx2}{TX2} sans TensorRT mentionné
    \citep[Table~2]{2019-ding-req-yolo-resource-aware-efficient} ;
  \item \glslink{tx2}{TX2} de 2017, framework \nr{}
    \citep[Table~IV]{2023-montgomerie-corcoran-satay-streaming-architecture-toolflow-accelerating} ;
  \item \glslink{jetson-nano}{Jetson Nano} reprise d'un autre article, avec un autre modèle
    \citep[Table~7]{2023-zhai-fpga-vehicle-detection-tracking-accelerator} ;
  \item \glslink{jetson-nano}{Jetson Nano} en FP32
    \citep[§5.3]{2024-zhang-end-end-acceleration-yolo-object} ;
  \item \glslink{xavier-nx}{Xavier NX} bridée à 10\,\glslink{u-w}{W}, puissance par jtop
    \citep[§IV, Table~VI]{arxiv-2407.17647-energy-efficient-artefact-detection-accelerator} ;
  \item \glslink{jetson-nano}{Jetson Nano} à batch~8 contre \glslink{zcu102}{ZCU102} à batch~1
    \citep[Table~VI]{2024-aggarwal-chameleon-dual-memory-replay-online} ;
  \item \glslink{agx-orin}{Jetson AGX Orin} exécutée en PyTorch à batch~1
    \citep[§VII-C]{arxiv-2604.21280-imagehd-energy-efficient-device-continual} ;
  \item aucune baseline GPU embarqué : Kim et al., Karki et al.\pp, \glslink{lp-yolo}{LP-YOLO}\pp, SMOF\pp.
\end{itemize}
Les protocoles équitables du corpus restent rares. fpgaConvNet face au Tegra~X1 (2017), à
budget de puissance égal, et \glslink{dac-sdc}{DAC-SDC} (2018), mesuré par les organisateurs, sont anciens.
Le seul protocole récent est celui de
\citet{arxiv-2304.11580-framework-benchmarking-real-time-embedded} (\glslink{zcu104}{ZCU104} contre AGX
Xavier, 2023), et la \glslink{jetson}{Jetson} y gagne en latence, en précision et en \glslink{fpsw}{FPS/W} \src{Table~1}.
Aucune source du corpus ne compare une FPGA embarquée à une \glslink{orin}{Jetson Orin}.

\subsection{Reproductibilité}

Aucune des 101 fiches ne mentionne de bitstream publié. Parmi la cinquantaine de fiches
expérimentales, une dizaine publient ou annoncent un code propre lié au FPGA :
Magalhães, FINN-R en partie, la branche hls4ml de l'\glslink{ae}{AE} du LHC, \glslink{req-yolo}{REQ-YOLO} (probablement la
partie entraînement), \glslink{dac-sdc}{DAC-SDC} (code des meilleures équipes), le streamer de FCMP,
l'extension HAWQ de Campos, SMOF\pp{} (GitHub,
\citealp{arxiv-2403.18921-smof-streaming-modern-cnns-fpgas}), \glslink{chameleon}{Chameleon} (GitHub annoncé,
\citealp{2024-aggarwal-chameleon-dual-memory-replay-online}) et \glslink{skynet}{SkyNet}\pp{} (lien annoncé,
absent du texte extrait). Les détecteurs les plus proches du cas d'application ne
publient rien : Schlosser et al., Karki et al.\pp, \glslink{lp-yolo}{LP-YOLO}\pp, ni Yan, Zhang 2024, Kim ou
Fata. Réserve : la lecture se limite au texte extrait, et les liens en note de bas de page
sont souvent perdus.
